% 附加材料：锚定运行时参数
% 2026-08-15 由主稿附录移出。主稿正文（9 页限制）不再包含本节，也不引用「附录」或「附加材料」。
% 用法：投稿时作为独立的 supplementary PDF 编译，或粘回主稿 \appendix 之后。
% 正文已内联的数值：三个可靠性阈值（0.2/0.45/0.5）、450 ms 与 1.0 s 有效性门槛、
% s=1.15、Delta_min=250 ms、状态年龄估计系数（升/降）0.5/0.05、StaticLock 入锁门槛族、死区、
% 锁内增益半衰期、累计漂移上限、头动放大上限 4x、w_C/w_D、epsilon，以及两个对照条件的参数。

\section{运行时参数}\label{sec:appendix}

表~\ref{tab:appendix-params}列出锚定运行时的完整参数取值。全部参数在实验一至实验三中预先固定，不按场景、物体或条件调整。

锁内更新所用的低增益按
$g_j = (1 - 2^{-\Delta t_j / h_{\mathrm{creep}}}) R_j
(1 - \rho_j^{\mathrm{head}})$
计算，其中$\rho_j^{\mathrm{head}} \in [0,1]$取头部线速度与角速度相对满容忍尺度（0.3~m/s与60~°/s）之比的较大者并截断至1，因此头动增强时更新增益相应减小。头动强度只放宽运动类判据（入锁速度门槛、明确运动、累计漂移上限与死区），放大系数为$1+3\rho_j^{\mathrm{head}}$（上限4倍）；可靠性判据不随之放松。头部由运动转为静止后设置600~ms沉降窗口，其间冻结运动类解锁判据，避免头动滑移尚未褪去时误解锁。

持续残差判据所用的带泄漏累积和按
\begin{equation}
E_{p,j} = 2^{-\Delta t_j / h_E} E_{p,j-1} + \frac{\Delta t_j}{t_{\mathrm{ref}}} R_j \max\!\left(0,\, d_{p,j} - d_{\mathrm{db}}\right)
\label{eq:cusum}
\end{equation}
更新，旋转通道同理；$h_E$为证据半衰期，$\Delta t_j$为相邻锁内观测的间隔，除以参考间隔$t_{\mathrm{ref}}$使累积速率与候选到达率无关。$E_{p,j}$或$E_{\theta,j}$超过上限即触发释放。观测共识为对接纳观测位姿的低增益指数平滑，其半衰期复用$h_E$。

两个对照条件的参数同样预先固定：One-Euro滤波的位置通道最小截止频率为0.8~Hz、速度系数为6、导数截止频率为2~Hz，旋转通道依次为1~Hz、1与2~Hz；预测式追踪（Smoothed KF Extrapolation）的预测时域上限为0.18~s，校正残差衰减半衰期为0.06~s。

\begin{table}[t]
\centering
\caption{锚定运行时参数。上段为观测准入与历史状态查询，中段为静止锚定，下段为有效性规则与VCD评分。}
\label{tab:appendix-params}
\small
\setlength{\tabcolsep}{3.4pt}
\begin{tabularx}{\columnwidth}{@{}>{\raggedright\arraybackslash}Xll@{}}
\toprule
参数 & 记号 & 取值 \\
\midrule
\multicolumn{3}{@{}l}{\emph{观测准入与历史状态查询}} \\
准入可靠性阈值 & $R_{\min}$ & 0.2 \\
卡尔曼过程噪声（平移/旋转） & --- & 0.002~m$^2$/s$^3$ / 0.2~rad$^2$/s$^3$ \\
卡尔曼量测噪声（平移/旋转） & --- & $4\times10^{-6}$~m$^2$ / $4\times10^{-4}$~rad$^2$ \\
查询滞后系数 & $s$ & 1.15 \\
最小查询滞后 & $\Delta_{\min}$ & 250~ms \\
状态年龄估计系数（升/降） & --- & 0.5 / 0.05 \\
\midrule
\multicolumn{3}{@{}l}{\emph{静止锚定（StaticLock）}} \\
入锁线速度门槛 & $v_{\mathrm{th}}$ & 50~mm/s \\
入锁角速度门槛 & $\omega_{\mathrm{th}}$ & 22~°/s \\
入锁可靠性门槛 & $R_{\mathrm{lock}}$ & 0.25 \\
入锁驻留时间 & $\tau_{\mathrm{dwell}}$ & 350~ms \\
入锁运动统计平滑系数 & --- & 0.5（逐观测） \\
死区（平移/旋转） & $d_{\mathrm{db}}$ / $\theta_{\mathrm{db}}$ & 8~mm / 3~° \\
锁内更新增益半衰期 & $h_{\mathrm{creep}}$ & 2.7~s \\
CUSUM 上限（平移/旋转） & $E_p^{\max}$ / $E_\theta^{\max}$ & 80~mm / 20~° \\
CUSUM 证据半衰期 & $h_E$ & 0.27~s \\
CUSUM 参考间隔 & $t_{\mathrm{ref}}$ & 0.2~s \\
累计漂移上限（平移/旋转） & --- & 15~mm / 12~° \\
速度逃逸倍率与时长 & --- & 2.5$\times$，400~ms \\
低可靠性释放（阈值/时长） & --- & 0.3，600~ms \\
重锁抑制时长 & --- & 1.0~s \\
接缝过渡逐帧衰减 & --- & 0.85 \\
头部沉降窗口 & --- & 600~ms \\
头动满容忍尺度（线/角） & --- & 0.3~m/s / 60~°/s \\
头动容忍放大上限 & --- & 4$\times$ \\
\midrule
\multicolumn{3}{@{}l}{\emph{有效性与 VCD 评分}} \\
追踪可靠性下限 & --- & 0.5 \\
缓冲保持时长 & --- & 无可靠观测 450~ms \\
判定丢失 & --- & 无可靠观测 1.0~s \\
重新获取触发（阈值/时长） & --- & 0.45，600~ms \\
重新获取冷却 & --- & 3.0~s \\
颜色/深度权重 & $w_C$ / $w_D$ & 0.2 / 0.8 \\
对数下限 & $\epsilon$ & 0.05 \\
深度容差（下限/距离比） & --- & 5~mm / 0.02 \\
内点率权重与残差倍率 & --- & 0.6 / 2.5 \\
有效深度覆盖率门槛 & --- & 0.10 \\
结构项权重上限与IQR门槛 & --- & 0.35 / 20~mm \\
LAB亮度通道权重 & --- & 0.3 \\
\bottomrule
\end{tabularx}
\end{table}
